% GENERATED FILE. Edit canonical lessons, then run npm run generate:books.
% canonical-source-hash: fnv1a64:3f9c08ae590339a0
% canonical-lessons: HI-C506-karigar, HI-C506-auzar, HI-C506-hunar, HI-C506-garhna, HI-C506-mistri

\chapter{The Craftsman's Bench}
\label{ch:hi-the-craftsman-s-bench}
\hlchaptermodality{506}

\begin{chapteropening}
\textbf{By the end of this chapter:} \emph{I can name skilled trades and the tools they use.}

Complete the last lesson of chapter 506: I can name skilled trades and the tools they use.
\end{chapteropening}

\section[\texorpdfstring{\hi{कारीगर} (kārīgar)}{कारीगर (kārīgar)}]{\hi{कारीगर} (kārīgar) — a craftsman; an artisan}

\label{lesson:HI-C506-karigar}



\begin{quote}
Before the new one: say the Hindi for turning over (in bed), then the Hindi for a stretch (of the body).
\end{quote}

\subsection*{You'll want to know: \hi{कारीगर}}
\textbf{\hi{कारीगर}} — \emph{kārīgar} — \textquotedblleft{}a craftsman; an artisan\textquotedblright{}.

\begin{quote}
\hi{यह} \hi{कुर्सी} \hi{गाँव} \hi{के} \hi{एक} \hi{कारीगर} \hi{ने} \hi{बनाई।}
\end{quote}

\begin{quote}
\emph{Yah kursī gā̃v ke ek kārīgar ne banāī.}
\end{quote}

\begin{quote}
A craftsman from the village made this chair.
\end{quote}

Masculine or feminine, by the person meant.

\begin{cousinweb}
\textbf{\hi{कारीगर}} (\emph{kārīgar}) is usually traced to Persian \emph{kārīgar}, \textquotedblleft{}a craftsman\textquotedblright{}.
\end{cousinweb}

\begin{grammarlens}[title={What you've built}]
One more word for describing what you lived and felt.
\end{grammarlens}

\subsection*{Your turn}
\begin{itemize}
\raggedright
  \item \emph{Say it:} \emph{kārīgar}
  \item \emph{Say it:} \emph{kārīgar}, once more
  \item \emph{Say it:} \emph{Yah kursī gā̃v ke ek kārīgar ne banāī.} — A craftsman from the village made this chair.
  \item \emph{Recall:} say the Hindi for turning over (in bed), then the Hindi for a stretch (of the body), then say \emph{kārīgar} again
\end{itemize}

\begin{tcolorbox}[breakable,colback=teal!4,colframe=teal!35!black,title={Before you move on}]
What does \hi{कारीगर} mean? (\textquotedblleft{}A craftsman; an artisan\textquotedblright{}.) Say it once more.
\end{tcolorbox}
\section[\texorpdfstring{\hi{औज़ार} (auzār)}{औज़ार (auzār)}]{\hi{औज़ार} (auzār) — a tool}

\label{lesson:HI-C506-auzar}



\begin{quote}
Before the new one: say the Hindi for a stretch (of the body), then the Hindi for a craftsman.
\end{quote}

\subsection*{You'll want to know: \hi{औज़ार}}
\textbf{\hi{औज़ार}} — \emph{auzār} — \textquotedblleft{}a tool\textquotedblright{}.

\begin{quote}
\hi{बढ़ई} \hi{अपने} \hi{औज़ार} \hi{लकड़ी} \hi{के} \hi{एक} \hi{डिब्बे} \hi{में} \hi{रखता} \hi{है।}
\end{quote}

\begin{quote}
\emph{Baṛhaī apne auzār lakṛī ke ek ḍibbe meṁ rakhtā hai.}
\end{quote}

\begin{quote}
The carpenter keeps his tools in a wooden box.
\end{quote}

Masculine.

\begin{cousinweb}
\textbf{\hi{औज़ार}} (\emph{auzār}) is usually traced to Arabic \emph{awzār}, a plural meaning \textquotedblleft{}loads\textquotedblright{} or \textquotedblleft{}gear\textquotedblright{}, which reached Hindi through Persian; the sense \textquotedblleft{}tools\textquotedblright{} is the one Persian and Hindi use.
\end{cousinweb}

\begin{grammarlens}[title={What you've built}]
One more word for describing what you lived and felt.
\end{grammarlens}

\subsection*{Your turn}
\begin{itemize}
\raggedright
  \item \emph{Say it:} \emph{auzār}
  \item \emph{Say it:} \emph{auzār}, once more
  \item \emph{Say it:} \emph{Baṛhaī apne auzār lakṛī ke ek ḍibbe meṁ rakhtā hai.} — The carpenter keeps his tools in a wooden box.
  \item \emph{Recall:} say the Hindi for a stretch (of the body), then the Hindi for a craftsman, then say \emph{auzār} again
\end{itemize}

\begin{tcolorbox}[breakable,colback=teal!4,colframe=teal!35!black,title={Before you move on}]
What does \hi{औज़ार} mean? (\textquotedblleft{}A tool\textquotedblright{}.) Say it once more.
\end{tcolorbox}
\section[\texorpdfstring{\hi{हुनर} (hunar)}{हुनर (hunar)}]{\hi{हुनर} (hunar) — a skill; a craft}

\label{lesson:HI-C506-hunar}



\begin{quote}
Before the new one: say the Hindi for a craftsman, then the Hindi for a tool.
\end{quote}

\subsection*{You'll want to know: \hi{हुनर}}
\textbf{\hi{हुनर}} — \emph{hunar} — \textquotedblleft{}a skill; a craft\textquotedblright{}.

\begin{quote}
\hi{मेरे} \hi{दादा} \hi{के} \hi{हाथों} \hi{में} \hi{बहुत} \hi{हुनर} \hi{था।}
\end{quote}

\begin{quote}
\emph{Mere dādā ke hāthoṁ meṁ bahut hunar thā.}
\end{quote}

\begin{quote}
My grandfather had great skill in his hands.
\end{quote}

Masculine.

\begin{cousinweb}
\textbf{\hi{हुनर}} (\emph{hunar}) is usually traced to Persian \emph{hunar}, \textquotedblleft{}skill, art\textquotedblright{}.
\end{cousinweb}

\begin{grammarlens}[title={What you've built}]
One more word for describing what you lived and felt.
\end{grammarlens}

\subsection*{Your turn}
\begin{itemize}
\raggedright
  \item \emph{Say it:} \emph{hunar}
  \item \emph{Say it:} \emph{hunar}, once more
  \item \emph{Say it:} \emph{Mere dādā ke hāthoṁ meṁ bahut hunar thā.} — My grandfather had great skill in his hands.
  \item \emph{Recall:} say the Hindi for a craftsman, then the Hindi for a tool, then say \emph{hunar} again
\end{itemize}

\begin{tcolorbox}[breakable,colback=teal!4,colframe=teal!35!black,title={Before you move on}]
What does \hi{हुनर} mean? (\textquotedblleft{}A skill; a craft\textquotedblright{}.) Say it once more.
\end{tcolorbox}
\section[\texorpdfstring{\hi{गढ़ना} (gaṛhnā)}{गढ़ना (gaṛhnā)}]{\hi{गढ़ना} (gaṛhnā) — to shape; to forge}

\label{lesson:HI-C506-garhna}



\begin{quote}
Before the new one: say the Hindi for a tool, then the Hindi for a skill.
\end{quote}

\subsection*{You'll want to know: \hi{गढ़ना}}
\textbf{\hi{गढ़ना}} — \emph{gaṛhnā} — \textquotedblleft{}to shape; to forge\textquotedblright{}.

\begin{quote}
\hi{लोहार} \hi{गरम} \hi{लोहे} \hi{से} \hi{एक} \hi{कुदाल} \hi{गढ़} \hi{रहा} \hi{है।}
\end{quote}

\begin{quote}
\emph{Lohār garam lohe se ek kudāl gaṛh rahā hai.}
\end{quote}

\begin{quote}
The blacksmith is shaping a hoe from hot iron.
\end{quote}

Past, with \hi{ने} (he or she …): \hi{उसने} \hi{लोहे} \hi{से} \hi{एक} \hi{बर्तन} \hi{गढ़ा}, \emph{usne lohe se ek bartan gaṛhā}.

\begin{cousinweb}
\textbf{\hi{गढ़ना}} (\emph{gaṛhnā}) is a tadbhava, inherited from \emph{*gaṭhati}, \textquotedblleft{}makes, forms\textquotedblright{}, a form CDIAL reconstructs from its descendants (CDIAL 3966). CDIAL gives the Hindi verb as \textquotedblleft{}to form by hammering\textquotedblright{}.
\end{cousinweb}

\begin{grammarlens}[title={What you've built}]
One more word for describing what you lived and felt.
\end{grammarlens}

\subsection*{Your turn}
\begin{itemize}
\raggedright
  \item \emph{Say it:} \emph{gaṛhnā}
  \item \emph{Say it:} \emph{gaṛhnā}, once more
  \item \emph{Say it:} \emph{Lohār garam lohe se ek kudāl gaṛh rahā hai.} — The blacksmith is shaping a hoe from hot iron.
  \item \emph{Recall:} say the Hindi for a tool, then the Hindi for a skill, then say \emph{gaṛhnā} again
\end{itemize}

\begin{tcolorbox}[breakable,colback=teal!4,colframe=teal!35!black,title={Before you move on}]
What does \hi{गढ़ना} mean? (\textquotedblleft{}To shape; to forge\textquotedblright{}.) Say it once more.
\end{tcolorbox}
\section[\texorpdfstring{\hi{मिस्त्री} (mistrī)}{मिस्त्री (mistrī)}]{\hi{मिस्त्री} (mistrī) — a mechanic; a mason}

\label{lesson:HI-C506-mistri}



\begin{quote}
Before the new one: say the Hindi for a skill, then the Hindi for to shape.
\end{quote}

\subsection*{You'll want to know: \hi{मिस्त्री}}
\textbf{\hi{मिस्त्री}} — \emph{mistrī} — \textquotedblleft{}a mechanic; a mason\textquotedblright{}.

\begin{quote}
\hi{हमारी} \hi{गाड़ी} \hi{बिगड़} \hi{गई}, \hi{इसलिए} \hi{हमने} \hi{एक} \hi{मिस्त्री} \hi{को} \hi{बुलाया।}
\end{quote}

\begin{quote}
\emph{Hamārī gāṛī bigaṛ gaī, isliye hamne ek mistrī ko bulāyā.}
\end{quote}

\begin{quote}
Our car broke down, so we called a mechanic.
\end{quote}

Masculine.

\begin{cousinweb}
\textbf{\hi{मिस्त्री}} (\emph{mistrī}) is usually traced to Portuguese \emph{mestre}, \textquotedblleft{}a master (of a trade)\textquotedblright{}.
\end{cousinweb}

\begin{grammarlens}[title={What you've built}]
One more word for describing what you lived and felt.
\end{grammarlens}

\subsection*{Your turn}
\begin{itemize}
\raggedright
  \item \emph{Say it:} \emph{mistrī}
  \item \emph{Say it:} \emph{mistrī}, once more
  \item \emph{Say it:} \emph{Hamārī gāṛī bigaṛ gaī, isliye hamne ek mistrī ko bulāyā.} — Our car broke down, so we called a mechanic.
  \item \emph{Recall:} say the Hindi for a skill, then the Hindi for to shape, then say \emph{mistrī} again
\end{itemize}

\begin{tcolorbox}[breakable,colback=teal!4,colframe=teal!35!black,title={Before you move on}]
What does \hi{मिस्त्री} mean? (\textquotedblleft{}A mechanic; a mason\textquotedblright{}.) Say it once more.
\end{tcolorbox}
